\documentclass[conference]{IEEEtran} \usepackage{url}
\usepackage[utf8]{inputenc}
\usepackage{url}
\usepackage{hyperref}
\usepackage{amsmath}
\usepackage{graphicx}

\title{Labwork 2}
\author{
	Le-Chi Thanh Lam\\
	\texttt{23BI14248}
}

\begin{document}

\maketitle

\section{Introduction}

This report details the implementation of a U-Net based segmentation model for
measuring fetal head circumference from ultrasound images as part of the second
labwork.

\section{Data Description}

The dataset is sourced from the HC18 Grand Challenge. It contains 1,334
two-dimensional ultrasound images of the fetal head standard plane. The
collection consists of 999 training images and 335 test images.

Each image has the resolution of $800 \times 540$ pixels with pixel size
ranging from 0.052 to 0.326 mm. Pixel sizes are provided via CSV files. For the
training set, manual annotations and head circumference measurements (in mm)
are included.

In this report, the original test set is not used due to the absence of public
labels. Instead, the 999 training images are split into 899 images for training
and 100 images for validation and performance measurement.

\section{Model architecture}

The model utilizes the U-Net architecture from the
\texttt{segmentation-models.pytorch}library \cite{smp}. A ResNet50 backbone
with ImageNet pretrained weights serves as the encoder. Since the input
images are grayscale, the model is configured with a single input channel.
For binary segmentation, the output class is set to 1. Spatial and channel
squeeze-and-exication (scSE) attention modules are integrated into the decoder.

\section{Training}

Images are resized to $512 \times 512$ pixels. Data augmentation includes random horizontal flips ($p=0.5$), shifting, scaling, and rotation. Pixel values are normalized and converted to tensors using \texttt{Albumentations} and \texttt{torchvision}.

The model is training for 40 epochs with a batch size of 8. We utilize the
Diceloss function and the AdamW optimizer with a learning rate of $1e-4$
and weight decay of $1e-5$. Mixed precision training is implemented via
\texttt{GradScaler}.

\section{Result}

Post-processing involves \texttt{cv2.findContours} to extract the mask boundary.
These contours are used to fit an ellipse via \texttt{cv2.fitEllipse}. The final
head circumference is calculated using Ramanujan's approximation for ellipse
circumference.

On the 100-image validation set, the model achieved an absolute mean difference
and standard deviation of $23.76 \pm 12.67$ mm. This performance is lower than
the top tier results on the HC18 leaderboard, indicating a need for further
hyperparameter tuning or architectural refinement.

\begin{thebibliography}{1}
\bibitem{smp}
P. Yakubovskiy, "Segmentation Models Pytorch," GitHub, 2020. [Online]. Available: \url{https://github.com/qubvel/segmentation_models.pytorch}.
\bibitem{hc18}
"HC18 - Grand Challenge," [Online]. Available: \url{https://hc18.grand-challenge.org/}.
\end{thebibliography}

\end{document}
